\fancyhead[C]{Section 13.2}
\fancyhead[R]{\dayfour}

\section*{\centering Chapter 13.2: Integrals of Vector Valued Functions}
\textbf{G2: Calculus of Curves.} I can compute tangent vectors to parametric curves and their velocity, speed, and acceleration. I can find equations of tangent lines to parametric curves. I can solve initial value problems for motion on parametric curves.

\subsection*{Mechanics}
\begin{enumerate}
    \item \question{Solve the following initial value problems: 
    \begin{enumerate}
        \item $\br''(t)=-\bi-\bj-\bk, \quad t\geq 0, \qquad \br'(0)=5 \bi, \qquad \br(0)=10\bi+10\bj+10\bk$
        \item $\br''(t)=(t,e^t,e^{-1}) \quad t\in \R, \qquad \br'(0)=(0,0,1), \qquad \br(0)=(0,1,1)$
    \end{enumerate}}
    {% answer goes here

    }
    {% solution goes here
    } 
\end{enumerate}
\subsection*{Applications}
\begin{enumerate}[resume]

	\item \question{A baseball is hit when it is $2.5$ ft above the ground. It leaves the bat with an initial velocity of $140$ ft/sec at a launch angle of $30^\circ$. At the instant the ball is hit, an instantaneous gust of wind blows against the ball, adding a component of $-14 \bi$ (ft/sec) to the ball's initial velocity. A $15$ ft high fence lies $400$ ft from the home plate in the direction of the flight. (Note that gravity, g $= -32 \bj$ ft/sec$^2$)
	
	\begin{enumerate}
		\item Find a vector equation for the path of the baseball.
		
		\item How high does the baseball go, and when does it reach maximum height?
		
		\item Find the range and flight time of the baseball, assuming that the ball is not caught.
		
		\item When is the baseball $20$ ft high? How far (ground distance) is the baseball from home plate at that height?
		
		\item Has the batter hit a home run? Explain.
	\end{enumerate}	}
    {% answer goes here

    }
    {% solution goes here
    } 
    \item \question{A fortified medieval city is surrounded walls with length 500m and height 15m. You are the commander of an attacking army and the closest you can get to the wall is 100m. Your plan is to set fire to the city by catapulting heated rocks over the wall at an initial speed
of 80 ms$^{-1}$. At what range of angles should you tell your army to set the catapult? You may assume the path of the rocks is perpendicular to the wall, and that you are as close as possible. \textit{[Note: in contrast to the previous problem, here we have $g = -9.8 \bj$ ms$^{-2}$].}}
    {% answer goes here

    }
    {% solution goes here
    } 
	
\end{enumerate}
\subsection*{Extensions}
\begin{enumerate}[resume]
    \item \question{Show that a projectile reaches three-quarters of its maximum height in half the time needed to reach its maximum height.}
    {% answer goes here

    }
    {% solution goes here
    } 
\end{enumerate}
{}

\iftoggle{answers}{
\fancyhead[R]{\dayfive}
\begin{center}{\large \textbf{Math 2551 Worksheet 5 Answers: Calculus of Vector-Valued Functions}}
\end{center}

\begin{enumerate}
	
	\item $\br(t)=\langle -\dfrac{1}{2}t^2+5t+10,-\dfrac{1}{2}t^2+10,-\dfrac{1}{2}t^2+10\rangle, \quad t\geq 0$.
	\item A baseball is hit when it is $2.5$ ft above the ground. It leaves the bat with an initial velocity of $140$ ft/sec at a launch angle of $30^\circ$. At the instant the ball is hit, an instantaneous gust of wind blows against the ball, adding a component of $-14 \hat{i}$ (ft/sec) to the ball's initial velocity. A $15$ ft high fence lies $400$ ft from the home plate in the direction of the flight. (Note that gravity, g $= 32$ ft/sec$^2$)
	
	\begin{enumerate}
		\item Sorry, no sketch. :)
		
		\item $\vect r(t) = (140\cos 30^\circ-14)t \hat{i}+ (2.5+(140\sin 30^\circ)t-16t^2)\hat{j} = (70\sqrt{3} -  14) t \hat{i} + (2.5+70t-16t^2) \hat{j}$.
		
		\item $y_{\text{max}}=\frac{(140 \sin 30^\circ)^2}{64}+2.5=\frac{70^{2}}{64}+2.5 = 79.0625$ ft., which is reached at $t = \frac{140 \sin 30^\circ}{32}=\frac{70}{32}=2.1875$ s.
		
		\item For the time, solve $y=2.5+70t-16t^2=0$ for $t$. Using quadratic formula, we have $t=4.41$s. Then, the range at $t=4.41$ is $x(4.41) = (140\cos 30^\circ-14)(4.41)=472.94$ ft. 
		
		\item For the time, solve $y=2.5+70t-16t^2=20$ for $t$. Using quadratic formula, we have $t=0.27, \ 4.11$ seconds. Then, the range at those times are $x(0.27) = 29$ ft and $x(4.11)=441$ ft.
		
		\item Yes, according to part (d), the ball is still 20 feet above the ground when it is 441 feet from home plate.
	\end{enumerate}	

\end{enumerate}
}{}
\iftoggle{solutions}
{
Solutions go here in the same format.
}{}
